\chapter{Zaključak}

\section{Dostignuća}

Razvijeni prototip sistema za udomljavanje životinja demonstrira sledeće:

\begin{enumerate}
    \item \textbf{Kompletan UML dizajn} - sve tražene dijagrame su kreirane
    \item \textbf{Jasna arhitektura} - troslojna struktura sa jasnom separacijom gladanja
    \item \textbf{Funkcionalan prototip} - sve funkcionalnosti za ocenu 10 su implementirane
    \item \textbf{Mogućnost demonstracije} - komponente komuniciraju međusobno
\end{enumerate}

\subsection{Funkcionalnosti dostignute za ocenu 10}

\begin{itemize}
    \item \checkmark Pregled svih korisnika sistema (administrator)
    \item \checkmark Pregled svih volontera iz udruženja (administrator udruženja)
    \item \checkmark Mogućnost importa i pregleda izvoda računa (administrator udruženja)
    \item \checkmark Sve prethodne funkcionalnosti (ocena 6-9)
\end{itemize}

\section{Primenjena načela}

\subsection{Dizajnerska načela}

\begin{itemize}
    \item \textbf{SOLID principi} - svaka klasa ima jasnu odgovornost
    \item \textbf{DRY (Don't Repeat Yourself)} - izbegavanje dupliranja koda
    \item \textbf{Separation of Concerns} - jasna separacija slojeva
    \item \textbf{Dependency Injection} - fleksibilna arhitektura
\end{itemize}

\subsection{UML i modeliranje}

\begin{itemize}
    \item Use case dijagram definiše sve interakcije
    \item Dijagram klasa prikazuje strukturu i relacije
    \item Sekvencijski dijagrami modeliraju glavne tokove
    \item Dijagrami aktivnosti pokazuju poslovne procese
\end{itemize}

\section{Mogućnosti proširenja}

\subsection{Kratkoročne mogućnosti}

\subsubsection{Notifikacioni sistem}
\begin{itemize}
    \item Implementacija push notifikacija
    \item Email notifikacije za važne događaje
    \item SMS notifikacije za udomitelje
\end{itemize}

\subsubsection{Napredne pretrage}
\begin{itemize}
    \item Filtriranje životinja po više kriterijuma
    \item Sortiranje po raznim poljima
    \item Tekstualna pretraga po opisu
\end{itemize}

\subsubsection{Izveštajni sistem}
\begin{itemize}
    \item Export izveštaja u PDF formatu
    \item Grafički prikazi statistike
    \item Automatsko slanje izveštaja menadžmentu
\end{itemize}

\subsection{Dugoročne mogućnosti}

\subsubsection{Mobilna aplikacija}
\begin{itemize}
    \item Native mobile app za iOS/Android
    \item Push notifikacije u realnom vremenu
    \item Offline režim za pregled životinja
\end{itemize}

\subsubsection{Integraciono sa društvenim mrežama}
\begin{itemize}
    \item Deljenje životinja na društvenim mrežama
    \item Ugradjena mapiranja na Facebook/Instagram
    \item Social login mogućnosti
\end{itemize}

\subsubsection{AI/ML mogućnosti}
\begin{itemize}
    \item Preporučivanje životinja na osnovu preferencija
    \item Automatska kategorizacija životinja iz slike
    \item Predictive analytics za potrebe donacija
\end{itemize}

\subsubsection{Sistem plaćanja}
\begin{itemize}
    \item Integracija sa Stripe ili PayPal
    \item Mogućnost online donacija
    \item Automatski izvršavanje plaćanja
\end{itemize}

\section{Zaključne napomene}

Razvijeni sistem predstavlja solidnu osnovu za upravljanje udomljavanjem životinja. Arhitektura je napravljena da bude skalabilna, održavana i lako proširiva za dodatne funkcionalnosti.

Korišćenje C\# i .NET platforme omogućava brz razvoj sa visokim nivoom sigurnosti i performansi. Avalonia framework je odličan izbor za multiplatformsku aplikaciju.

UML dijagrami pružaju jasan uvid u dizajn sistema i mogu služiti kao referenca za budući razvoj i održavanje.

\section{Literatura i resursi}

\begin{itemize}
    \item Specifikacija projekta iz predmeta SIMS
    \item .NET dokumentacija: https://docs.microsoft.com/dotnet
    \item Avalonia dokumentacija: https://docs.avaloniaui.net
    \item PostgreSQL dokumentacija: https://www.postgresql.org/docs
    \item PlantUML dokumentacija: https://plantuml.com
\end{itemize}

\vspace{2cm}


